%========================================================================
\section{Contexte \& fil rouge}
% ⏱ ~4 min  |  Objectif : poser le décor, montrer l'étendue du rôle, ENTONNOIR vers la question phare
%========================================================================

\begin{frame}{Ecocea La Fabrique \& le projet Polène}
  % 🎤 ~1 min. Qui est l'agence, sur quoi j'étais. Rester factuel et rapide.
  \begin{columns}[T]
    \column{.55\textwidth}
    \begin{itemize}
      \item Intégrateur e-commerce ($\sim$60 ingénieurs), écosystème \alert{Shopify} \& SFCC
      \item Comptes retail haut de gamme : Polène, Céline, Jacquemus
      \item Mon terrain : \alert{Polène} --- déployé dans $>$20 pays, livré en continu
    \end{itemize}
    \column{.45\textwidth}
    \placeholder{3.2cm}{logo Polène / capture polene-paris.com}
  \end{columns}
\end{frame}

\begin{frame}{Un rôle qui s'est élargi en 3 temps}
  % 🎤 ~1 min 30. C'est LA slide qui annonce la largeur. Montrer la timeline, dire que je vais ZOOMER sur le 1er bloc.
  % 👉 A remplacer par une frise visuelle (TikZ ou image). Stub ci-dessous.
  \placeholder{3.4cm}{frise : Delivery QA  $\to$  Product / PO  $\to$  Discovery GTT (LNG)}
  \vskip 4pt
  \footnotesize
  \begin{columns}[T]
    \column{.33\textwidth}\textbf{Jan–Mar --- QA / Delivery}\\ Homologation, TNR, MEP multi-pays
    \column{.33\textwidth}\textbf{Avr–Mai --- Product}\\ Rédaction de tickets, audits UX/UI, appels d'offres
    \column{.33\textwidth}\textbf{Mai–Juin --- GTT / LNG}\\ Benchmark, discovery, ateliers
  \end{columns}
  \vskip 4pt
  % 🎤 « Je vais zoomer sur le premier bloc --- le projet le plus structurant --- puis montrer que la démarche se rejoue ailleurs. »
\end{frame}

\begin{frame}{Le fil rouge de cette soutenance}
  % 🎤 ~1 min 30. Poser le fil rouge transverse PUIS la problématique précise du projet phare.
  \begin{block}{Un même geste d'ingénieur, répété sur 3 missions}
    Transformer une réalité \alert{non structurée, non documentée, mouvante} en un livrable
    \alert{reproductible et transmissible} --- l'IA comme levier, mon jugement comme garde-fou.
  \end{block}
  \vskip 6pt
  \begin{alertblock}{Problématique du projet phare (QA)}
    Comment garantir la qualité du Delivery d'un site e-commerce multi-pays en exploitation,
    avec un dispositif opérationnel restreint et des cadences soutenues ?
  \end{alertblock}
\end{frame}
